\documentclass[12pt,a4paper]{article}
\usepackage[utf8]{inputenc}
\usepackage[T1]{fontenc}
\usepackage[spanish]{babel}
\usepackage{geometry}
\usepackage{hyperref}
\usepackage{graphicx}
\usepackage{float}
\usepackage{caption}
\geometry{margin=2.2cm}

\title{Documentacion Final --- Simulador de Gestion de Procesos Mini PC\\\large Version LaTeX con UML en imagenes}
\author{Proyecto SimuladorMiniPC}
\date{\today}

\begin{document}
\maketitle
\tableofcontents
\newpage

\section{Descripcion del Proyecto}
Version de exportacion LaTeX orientada a lectura tecnica, con los diagramas UML incorporados como imagenes PNG.

\section{Inventario de Paquetes}
La arquitectura se organiza en 14 paquetes, con el kernel como orquestador del ciclo de ejecucion y subsistemas especializados para CPU, memoria, E/S, interrupciones, almacenamiento y metricas.

\section{Clases Clave --- UML en Imagenes}
Esta version reemplaza las vistas en texto/markdown por figuras insertadas, de forma equivalente a un documento HTML con imagenes embebidas.

\subsection{Diagrama UML de Clases}
\begin{figure}[H]
  \centering
  \includegraphics[width=\textwidth]{diagramas/diagrama_clases_completo.png}
  \caption{Diagrama UML de clases del simulador (vista integral).}
\end{figure}

\subsection{Diagrama UML de Paquetes}
\begin{figure}[H]
  \centering
  \includegraphics[width=\textwidth]{diagramas/diagrama_paquetes.png}
  \caption{Diagrama UML de paquetes y dependencias principales.}
\end{figure}

\subsection{Diagrama UML de Secuencia de Ejecucion}
\begin{figure}[H]
  \centering
  \includegraphics[width=\textwidth]{diagramas/diagrama_secuencia_ejecucion.png}
  \caption{Secuencia de ejecucion por tick entre kernel, scheduler, CPU y gestores.}
\end{figure}

\subsection{Diagrama UML de Flujo de Excepciones}
\begin{figure}[H]
  \centering
  \includegraphics[width=0.95\textwidth]{diagramas/diagrama_flujo_excepciones.png}
  \caption{Flujo UML de manejo de errores y excepciones de simulacion.}
\end{figure}

\section{Pipeline de Ejecucion}
El pipeline mantiene diez etapas fijas por tick: reloj, admision, avance de E/S, planificacion, ejecucion, manejo de eventos, metricas, notificacion y cierre.

\section{Gestion de Memoria}
Modelo first-fit sobre RAM fisica y traduccion base+offset para direccionamiento virtual.

\section{Interrupciones y E/S}
Soporte para INT 20H, INT 10H, INT 09H e INT 21H con latencias por dispositivo.

\section{Manejo de Errores}
Se mantienen rutas de parse y runtime con contencion del fallo por proceso.

\section{Metricas de Rendimiento}
Utilizacion de CPU, throughput, tiempos promedio y ratio de retorno.

\section{Conjunto de Instrucciones ASM}
Conjunto de instrucciones con peso de CPU y semantica operacional.

\section{Estados del Proceso}
Transiciones NEW, READY, RUNNING, BLOCKED y TERMINATED con reglas de negocio del estado.

\section{Manual de Usuario}
Operacion guiada de la interfaz, lectura de paneles y resolucion de incidencias comunes.

\section{Manual Tecnico}
Arquitectura, configuracion externa y lineamientos de mantenimiento.

\section{Objetivos Alcanzados}
Cumplimiento de objetivos academicos de modelado de SO, planificacion, memoria, E/S y metricas.

\section{Bibliografia}
Referencias de sistemas operativos, patrones y calidad de software.

\end{document}
